% A.58 full independent root and DE proof/source/certificate audits PASS.
\subsection{The actual five-adic Heisenberg model}

\begin{statement}{A.58}{Proved}{All-prime fixed groups in the genus-fifty-one Heisenberg cover}
Let \(C\), \(q:C\to E\), and \(Y=C\setminus q^{-1}(E[5])\)
be the actual degree-one-hundred-twenty-five Heisenberg cover of A.54:
\[
              E:y^2=x^3-x,\qquad z^5=\psi _5(x).
\]
After a finite extension at five its proper stable reduction has
one ordinary elliptic component, five genus-six components of the
form \(V^5-V=U^4\), and ten genus-two components of the form
\(V^5-V=U^2\), over the algebraic residue field. Every other
component is rational, and the dual graph is a tree. The elliptic
component maps radicially with degree five to \(E\); each other
positive-genus component maps to a distinct point outside the
reduced support of \(E[5]\).

Every selected cusp of \(Y\) has a second cusp joined to it by
a rational path in the pointed stable fibre. Consequently, for
every coefficient prime, every selected cusp packet, and every
indicated open derived subgroup, the full geometric pro-prime
fixed subgroup is exactly the selected cusp inertia. Together
with A.54--A.55, both assertions of Task A.5 are therefore proved
at every prime for this actual Heisenberg cover.
\end{statement}

\begin{proof}
Put
\[
\begin{aligned}
 g(x)&=3x^2+1,\\
 h(x)&=\frac{\psi _5(x)-g(x)^5}{5}\\
     &=x^{12}-61x^{10}-102x^8+6x^6-43x^4+7x^2.
\end{aligned}                                                   \tag{A.58.1}
\]
These are exact polynomial identities in the actual covering
field. At five the indicated elliptic model is smooth and ordinary.
For example its Cartier coefficient is the coefficient of \(x^4\)
in \((x^3-x)^2\), namely \(-2\ne0\pmod5\).
We work over a finite local extension containing the constants
used below and admitting semistable reduction; such constants
may also be included in a finite global field of definition.

\paragraph{The actual elliptic main component.}
Choose \(\pi^5=5\), and set \(z=g(x)+\pi u\). The exact equation is
\[
 u^5+\pi g^4u+2\pi^2g^3u^2+2\pi^3g^2u^3+\pi^4gu^4=h(x),
                   \qquad y^2=x^3-x.                 \tag{A.58.2}
\]
Its special function field is \(k(E)(h^{1/5})\). In
characteristic five the exact factorization is
\[
          dh=2x(x^2-2)^3(x^2-x+1)(x^2+x+1)\,dx\ne0.
                                                        \tag{A.58.3}
\]
The one-variable function field \(k(E)\) has \(p\)-basis
of cardinality one. Hence this extension is \(k(E)^{1/5}\),
and its smooth projective normalization is the Frobenius twist
of \(E\), of genus one. The chart is smooth at its special
generic point: the nonzero differential in (A.58.3) gives the
required Jacobian rank together with the elliptic equation.
Thus this is an actual positive-genus component of a proper
model of \(C\), with degree-five radicial map to \(E\).

\paragraph{The fifteen residue centres.}
On the open subset where \(g\) is a unit put
\[
              F=\frac{\psi _5}{g^5}=1+5H,
                         \qquad H=\frac h{g^5}.        \tag{A.58.4}
\]
In characteristic five, \(dH=dh/g^5\). Its zeros, counted
on the elliptic curve with a genuine local parameter, are as follows:
\[
\begin{array}{c|c|c|c}
 \text{centres}&\text{number of points}&\operatorname{ord}(dH)&k\\
 \hline
 (0,0)&1&3&4\\
 x^2=2&4&3&4\\
 (1,0),\ (-1,0)&2&1&2\\
 x^4+x^2+1=0&8&1&2
\end{array}                                                   \tag{A.58.5}
\]
At a two-torsion point \((a,0)\), the parameter is \(y\),
and \(x-a\) has order two while \(dx\) has order one.
At all other displayed points \(x-a\) is a parameter.
The polynomial factors in (A.58.3) are square-free and
pairwise disjoint except for the displayed multiplicity three;
they meet neither \(g=0\) nor any additional two-torsion point.
This proves (A.58.5) directly. Its zero orders total twenty-five.

Choose lifts of these centres in a finite \emph{unramified}
extension \(R_0/\Z_5\). For the non-two-torsion centres take
\(a^2=2\) or \(a^4+a^2+1=0\), respectively, and a root
\(y_0^2=a^3-a\). All required residue roots are simple,
and \(y_0\) is a unit. The other lifts are the actual
points \((0,0),(1,0),(-1,0)\). The unramified choice
ensures that a coefficient reducing to zero is divisible
by \(5\), not merely by a smaller uniformizer after a
ramified extension.

\paragraph{Exact algebraic tail charts.}
For any selected lift \((a,y_0)\), let
\[
 F_0=F(a),\qquad H_0=H(a),\qquad
 N_a(x)=h(x)-H_0g(x)^5,
 \qquad z_0^5=F_0.                                    \tag{A.58.6}
\]
Here \(F_0\in R_0^\times\), \(F_0\equiv1\pmod5\),
and \(N_a(a)=0\) exactly. Enlarge the field to choose
\(\eta^4=-5\) and \(\rho^k=\eta\), where \(k\) is
the last column of (A.58.5). The exact covering-coordinate change is
\[
                    z=g(x)z_0(1+\eta V).
\]
After subtraction of the constant term and division by \(5\eta\),
the second equation is
\[
 g(x)^5F_0
   \bigl(V+2\eta V^2+2\eta^2V^3+\eta^3V^4-V^5\bigr)
                         =\frac{N_a(x)}{\eta}.        \tag{A.58.7}
\]
The equality follows by the binomial theorem and \(\eta^4=-5\);
no formal truncation has been made.

For \(y_0\ne0\), substitute
\(x=a+\rho U,\ y=y_0+\rho W\). The first exact equation is
\[
 2y_0W+\rho W^2
       =(3a^2-1)U+3a\rho U^2+\rho^2U^3.               \tag{A.58.8}
\]
Write \(N_a(a+t)=\sum_{j\ge1}n_jt^j\). The order
in (A.58.5), and \(k<5\), imply
\[
 n_j\in5R_0\ (1\le j<k),\qquad n_k\in R_0^\times.
                                                        \tag{A.58.9}
\]
Indeed \(N_a=g^5(H-H_0)\), and in characteristic five
the first nonconstant term of \(H\) has degree \(k\).
Thus \(N_a(a+\rho U)/\eta\) is integral. For \(j<k\)
its coefficient has valuation at least
\(1+(j-k)/(4k)>0\); for \(j>k\) it has valuation
at least \((j-k)/(4k)>0\). Its reduction is exactly
\(\overline n_k U^k\). Equation (A.58.8) eliminates
\(W\) linearly on the special fibre.

For \(y_0=0\), put \(\delta=3a^2-1\), a unit, and use
\(y=\rho U,\ x=a+\rho^2W\). The first equation now is
\[
           U^2=\delta W+3a\rho^2W^2+\rho^4W^3.        \tag{A.58.10}
\]
The first nonzero term of \(\overline N_a(a+t)\) has
degree \(k/2\), with nonzero coefficient. All lower
coefficients are divisible by \(5\). Substitution in
(A.58.7), with \(\rho^k=\eta\), is again integral;
its reduction is a nonzero multiple of \(W^{k/2}\).
Equation (A.58.10) eliminates
\(W=\overline\delta^{-1}U^2\).

In both cases the resulting special equation is therefore
\[
                         V-V^5=c_aU^k,
                         \qquad c_a\in k^\times.      \tag{A.58.11}
\]
The two equation partials in \(W,V\) are units on the
special fibre, after inverting \(g(x)\), whose reduction
there is a nonzero constant. These are smooth algebraic
charts over the enlarged valuation ring. Their generic
coordinate changes are invertible on the actual curve.
The special affine curve is geometrically integral: a
polynomial of degree \(k=2\) or \(4\) is not an
Artin--Schreier coboundary, as its pole order is not
divisible by five. Its smooth completion has one point
over infinity. To compute the ramification directly, put
\(t=U^{-1}\), so the extension valuation gives
\(v(t)=5\) and \(v(V)=-k\). The element
\(tV^{4/k}\) is a uniformizer. Under a nontrivial
translation \(V\mapsto V+b\) its difference has
valuation \(k+1\), since \(4/k\) is a unit modulo
five. The lower jump is therefore \(k\), the different
has degree \(4(k+1)\), and Hurwitz gives genus
\((5-1)(k-1)/2=2(k-1)\). Over the algebraic residue
field rescaling \(U\) gives the models in the statement.

\paragraph{Persistence and exhaustion of the proper genus.}
Compactify these algebraic charts and take regular
resolutions. Each smooth special generic point defines
a component with the displayed positive genus. The
factorization of birational maps of regular arithmetic
surfaces into point blowups preserves every positive-genus
component on passage to the stable model; see
\cite[Theorem 9.2.2 and Remark 9.2.6, pp.395--396,
Theorem 9.3.21, p.422]{fghLiu}, as used explicitly in A.52.
The fifteen tail valuations are distinct because their
restrictions to the good elliptic model have the fifteen
different centres (A.58.5). The main valuation lies over
its special generic point, so is distinct from all of them.

Their genera already sum to
\[
                         1+5\cdot6+10\cdot2=51=g(C).
                                                        \tag{A.58.12}
\]
The genus formula for a connected semistable fibre forces
every further component to be rational and the dual graph
to have first Betti number zero. Thus the reduction is of
compact type. Each Artin--Schreier tail has \(p\)-rank
zero by the one-branch-point Deuring--Shafarevich formula
\cite[Corollary 1.8, pp.35--36]{fghCrew1984};
the main elliptic component is ordinary. In particular the
good special Jacobian has five-rank one, although this
rank calculation is not needed for the fixed-group criterion.
No unramified-cover version of Raynaud's theorem is
applied to this ramified degree-five morphism.

\paragraph{The actual map and the cusp clusters.}
The proper map \(q\) extends from a regular semistable
model of \(C\) to the good elliptic model. Resolve its
indeterminacy, and then descend through the resulting
point blowups: their exceptional curves are rational,
so every map from one of them to the elliptic special
fibre is constant. The same applies to stable contraction;
compare \cite[Theorem 9.2.7, pp.396--397]{fghLiu} and
the detailed argument in A.52. The main component has
degree five by (A.58.2). Pulling back a degree-one line
bundle on the target gives total degree five on every
fibre, so all other components map constantly. Each
connected subtree off the main component attaches at one
point, and maps to the image of that point. The main
map is radicial, so a given reduced point of \(E\)
has a unique possible attachment point on that component.

On the ordinary elliptic curve the twenty-five geometric
generic points of \(E[5]\) specialize five at a time
to its five geometric special points, by the connected--etale
sequence for its ordinary five-divisible group
\cite[Sections 1.4 and 2.2]{fghTate1967Pdivisible}.
In these coordinates
their reduced support is precisely
\[
                    \{O\}\ \cup\ \{x^2=3\}.          \tag{A.58.13}
\]
The finite support follows also from
\(\overline\psi _5=(3x^2+1)^5\); its four finite
points together with \(O\) exhaust the ordinary
special five-torsion. This support is disjoint from
every centre in (A.58.5).

The map \(q\) is totally ramified over each generic
point of \(E[5]\), with one cusp in each fibre, by
the division-polynomial divisor of A.54. Its twenty-five
cusps consequently form five clusters of five marks over
the five points (A.58.13). None can lie on a positive-genus
tail or on its connecting subtree, since those subtrees
map to one of the different points (A.58.5).
Every cluster therefore lies over a single main-component
point with only rational subtrees above it.

Separate the marked sections on a semistable model. Since
each such cluster has at least two marks, its main-component
point must be replaced by a rational marked tree. No mark
can remain at its attachment point on the main component,
because that point is a node in the separated marked model.
All joining paths within that cluster consequently lie
on rational components, including their endpoints.
Stable marked contractions preserve this property.

For each chosen cusp select another in the same cluster.
The precise pointed hypothesis of A.41--A.42 now holds.
Those statements give the full integral fixed-inertia
theorem for every coefficient prime, every packet, and
every permitted open derived subgroup. The all-prime
normalizer theorem was already proved through the actual
quotient maps in A.54--A.55. This completes both assertions
for the specified Heisenberg cover, without a cofinality
or fixed-lifting assertion for other covers.
\end{proof}
